\chapter{SIMPULAN DAN SARAN}
\label{ch:simpulan}

% =====================================================================
\section{Simpulan}
\label{sec:simpulan}
% =====================================================================

Penelitian ini telah membangun prototipe sistem pelaporan masalah lingkungan berbasis fusi multimodal citra dan teks yang siap diintegrasikan dengan ekosistem aplikasi pemerintah daerah. Berdasarkan rumusan masalah yang dipaparkan pada Bab~\ref{ch:pendahuluan}, simpulan penelitian ini adalah sebagai berikut.

\begin{enumerate}
    \item \textbf{Arsitektur klien-server berbasis Flutter, FastAPI, dan Supabase.} Sistem prototipe SmartCityApps berhasil diimplementasikan dengan klien Android berbasis Flutter, server inferensi FastAPI yang memuat tiga \textit{artifact} ONNX, dan \textit{backend} Supabase yang menyediakan otentikasi, basis data PostgreSQL, dan penyimpanan berkas. Tiga peran pengguna (Warga Pelapor, Operator Instansi, dan Moderator) berhasil dilayani dengan akses yang dibatasi melalui Row Level Security pada tabel \texttt{profiles}, \texttt{reports}, dan \texttt{report\_history}. Konsumsi sumber daya pada perangkat pengguna tetap minimum karena seluruh komputasi inferensi dipusatkan di server.

    \item \textbf{Kombinasi DINOv3 \textit{large}, \textit{multilingual-E5 large}, dan CatBoost dengan PGS.} Pada set uji yang berisi 9.266 sampel Jakarta CRM, konfigurasi terbaik mencapai \textit{akurasi} sebesar 0,8074, \textit{macro F1 score} sebesar 0,7747, dan \textit{balanced accuracy} sebesar 0,7916, mengungguli baseline modalitas teks saja (\textit{macro F1} 0,7515) sebesar 2,32 \textit{percentage point} dan baseline modalitas citra saja (\textit{macro F1} 0,5973) sebesar 17,74 \textit{percentage point}. Hasil ini menunjukkan bahwa skema fusi awal dengan \textit{classifier head} CatBoost yang dikalibrasi PGS efektif untuk klasifikasi laporan multimodal Bahasa Indonesia.

    \item \textbf{Paritas numerik PyTorch terhadap ONNX.} Validasi paritas pada seluruh komponen \textit{pipeline} (image encoder, text encoder, dan classifier head) memberikan selisih probabilitas absolut maksimum di bawah ambang $10^{-5}$, yang memastikan prediksi top-1 identik antara jalur pelatihan PyTorch dan jalur produksi ONNX. Server FastAPI dapat dipakai sebagai jalur inferensi produksi tanpa risiko divergensi numerik dengan jalur pelatihan.

    \item \textbf{Routing instansi terdekat.} Backend berhasil menggabungkan keluaran model 9 instansi dengan koordinat latitude-longitude laporan untuk memilih instansi tujuan terdekat yang relevan. Hasil routing yang disimpan mencakup identitas instansi, nama instansi, kategori instansi, estimasi jarak, dan metode routing, sehingga laporan tidak hanya diklasifikasikan, tetapi juga diarahkan ke unit penanganan yang paling dekat dengan lokasi kejadian.

    \item \textbf{Pemilihan inferensi \textit{on-server} terhadap \textit{on-device}.} Skema inferensi \textit{on-server} berbasis FastAPI dipilih karena total ukuran \textit{artifact} ONNX, dengan komponen \textit{image encoder} DINOv3 \textit{large} saja sebesar 1,15 GB, jauh melampaui batas wajar paket aplikasi Android. Konsekuensi dari pemilihan ini adalah ketergantungan pada ketersediaan jaringan dan pada layanan FastAPI sebagai titik inferensi tunggal, sekaligus keuntungan berupa pembaruan model yang terpusat tanpa harus merilis ulang aplikasi pada \textit{Play Store}, konsumsi daya yang rendah pada perangkat pengguna, dan kemampuan menjalankan model berskala besar yang tidak praktis untuk \textit{deployment} sepenuhnya \textit{on-device}. Skema \textit{on-server} ini mewarisi tiga keterbatasan yang perlu dicatat. Pertama, server inferensi pada prototipe ini di-\textit{deploy} pada HuggingFace Spaces, yaitu platform yang ditujukan untuk demonstrasi dan prototipe, bukan untuk lingkungan \textit{production-grade}: platform tersebut tidak memberikan jaminan ketersediaan layanan (\textit{Service Level Agreement}), berpotensi mengalami \textit{cold-start} setelah periode tidak aktif, dan menyediakan sumber daya komputasi yang terbatas. Konfigurasi ini memadai untuk pembuktian konsep pada penelitian ini, tetapi sebelum operasionalisasi nyata pada layanan publik, server perlu dipindahkan ke infrastruktur produksi terdedikasi. Kedua, terdapat ketergantungan pada ekosistem HuggingFace pada tataran pustaka dan lisensi: \textit{pipeline} serving masih mengimpor \texttt{transformers} dan \texttt{huggingface\_hub} untuk tokenisasi teks, sedangkan kedua \textit{encoder} beku (DINOv3 \textit{large} dan \textit{multilingual-E5 large}) berasal dari HuggingFace Hub beserta lisensi hulunya. Bobot model memang sudah diekspor ke berkas ONNX lokal sehingga inferensi tidak memerlukan akses jaringan ke Hub saat permintaan berlangsung; namun reproduktibilitas proses ekspor dan kepatuhan lisensi, khususnya ketentuan penggunaan DINOv3, tetap perlu ditinjau sebelum penggunaan produksi. Ketiga, pemusatan komputasi di server membuat biaya inferensi (kebutuhan GPU atau CPU berkapasitas tinggi untuk memuat dua \textit{encoder} besar) bersifat berkelanjutan dan meningkat seiring volume laporan, sehingga efisiensi biaya menjadi pertimbangan kelayakan yang belum dioptimasi pada prototipe ini.

    \item \textbf{Usabilitas tergolong ``dapat diterima'' (SUS rata-rata 71,6).} Evaluasi pengalaman pengguna menggunakan \textit{System Usability Scale} (SUS) terhadap 30 responden Warga Pelapor menghasilkan skor rata-rata 71,6 (median 72,2; simpangan baku 14,6; interval kepercayaan 95\% [66,1; 77,0]). Skor ini berada di atas ambang rata-rata industri 68 dan tergolong kategori \textit{dapat diterima} dengan deskriptor kata sifat \textit{good}, sedangkan konsistensi internal instrumen memadai (\textit{Cronbach's} $\alpha = 0{,}74$). Analisis per-pernyataan menunjukkan fungsionalitas inti dinilai kuat, sementara kebutuhan pembiasaan awal (\textit{learnability}) dan konsistensi antarmuka menjadi aspek terlemah. Perlu dicatat bahwa instrumen yang digunakan merupakan adaptasi sembilan butir dari SUS sepuluh butir baku (satu pernyataan negatif tidak disertakan dan skor dinormalisasi dengan faktor $100/36$), sehingga skor tidak sepenuhnya setara dengan SUS baku dan ukuran sampel yang terbatas membuat klaim usabilitas bersifat indikatif. Selain itu, SUS hanya mengukur kepuasan yang dipersepsikan sehingga efektivitas dan efisiensi tugas (\textit{task success rate}, \textit{time-on-task}, dan \textit{error rate}) tidak terukur, dan evaluasi terbatas pada peran Warga Pelapor sedangkan peran Operator Instansi dan Moderator belum diujikan kepada pengguna; ketiga hal ini menjadi keterbatasan utama evaluasi usabilitas pada penelitian ini.
\end{enumerate}

% =====================================================================
\section{Saran}
\label{sec:saran}
% =====================================================================

Berdasarkan temuan dan keterbatasan penelitian ini, beberapa arah penelitian lanjutan diusulkan sebagai berikut.

\begin{enumerate}
    \item \textbf{Distilasi dan kuantisasi untuk varian \textit{on-device}.} Penelitian lanjutan dapat menelaah distilasi pengetahuan dari DINOv3 \textit{large} ke arsitektur ringan seperti DINOv3 ViT-Small atau DINOv3 ViT-Base, dipadu dengan kuantisasi INT8 menggunakan ONNX Runtime Mobile. Varian ringan ini dapat dipakai sebagai jalur \textit{fallback} ketika koneksi jaringan tidak tersedia, sehingga warga pelapor tetap dapat memperoleh prediksi awal walaupun dengan akurasi yang sedikit lebih rendah.

    \item \textbf{Penambahan modalitas tambahan.} Dataset pelaporan publik dapat diperkaya dengan modalitas tambahan, contohnya audio (suara latar dari lokasi kejadian) yang berguna untuk laporan kebisingan, geolokasi yang dipakai sebagai fitur eksplisit pada vektor fusi sehingga model dapat memanfaatkan pola spasial, dan metadata waktu pelaporan untuk menangkap pola harian atau musiman.

    \item \textbf{Active learning dengan koreksi manual pengguna.} Mekanisme koreksi kategori secara manual oleh pengguna pada layar tinjauan AI dapat diubah menjadi sumber sinyal label baru. Setiap koreksi yang konsisten lintas pengguna dapat dipakai untuk melakukan pelatihan ulang model secara periodik, sehingga akurasi model meningkat dari waktu ke waktu seiring penggunaan aplikasi.

    \item \textbf{Perbaikan antarmuka berdasarkan temuan usabilitas.} Analisis per-pernyataan SUS dan umpan balik kualitatif menunjukkan bahwa hambatan utama terletak pada kurva pembelajaran awal, konsistensi antarmuka, dan kejelasan umpan balik setelah laporan dikirim. Penyempurnaan yang diusulkan meliputi penambahan \textit{onboarding} atau panduan singkat saat pertama kali memakai aplikasi untuk menurunkan kebutuhan pembiasaan, penyeragaman komponen dan istilah antarmuka lintas layar agar lebih konsisten, serta penguatan konfirmasi dan pelacakan status pasca-pengiriman laporan (mis. notifikasi dan linimasa status) agar pengguna memperoleh kepastian bahwa laporannya diproses. Evaluasi usabilitas ulang setelah perbaikan, idealnya dengan SUS sepuluh butir baku dan responden yang lebih beragam, dapat memverifikasi peningkatan skor.

    \item \textbf{Integrasi dengan API resmi pemerintah.} Penelitian lanjutan dapat melakukan integrasi langsung dengan API resmi platform LAPOR! SP4N yang dikelola oleh KemenPANRB atau platform JAKI yang dikelola oleh Pemerintah Provinsi DKI Jakarta. Dengan integrasi tersebut, hasil triase otomatis dapat langsung diteruskan ke kanal pemerintah, sehingga prototipe ini dapat berfungsi sebagai \textit{front-end} pelaporan multimodal yang melengkapi sistem yang sudah ada.

    \item \textbf{Registrasi lokasi instansi dan perutean berbasis wilayah administratif.} Koordinat instansi pada prototipe ini diperoleh melalui \textit{geocoding} alamat kantor resmi, tetapi cakupannya belum lengkap (enam belas kantor representatif, dengan suku dinas dan entri tingkat kelurahan diproksikan ke kompleks Kantor Walikota kotamadya). Penelitian lanjutan dapat menambahkan fitur registrasi lokasi mandiri oleh masing-masing Operator Instansi yang tersimpan pada Supabase, sehingga koordinat tujuan bersifat otoritatif dan mutakhir. Selain itu, perutean berbasis jarak Haversine dapat diganti dengan perutean berbasis batas wilayah administratif (kelurahan atau kecamatan) agar penugasan laporan sesuai dengan pembagian kewenangan instansi yang sebenarnya, bukan semata kantor yang terdekat secara garis lurus.

    \item \textbf{Perluasan dataset lintas daerah.} Dataset Jakarta CRM dapat diperluas dengan menambahkan dataset dari platform LaporGub Jawa Tengah \citep{zahroFairerPublicComplaint2025} atau dari pemerintah daerah lain di Indonesia. Perluasan ini memberikan kesempatan untuk menguji generalisasi model terhadap pola pelaporan di daerah dengan struktur instansi yang berbeda.

    \item \textbf{Penanganan kelas minoritas.} Tiga kelas minoritas pada Jakarta CRM (Dinas Cipta Karya, Tata Ruang, dan Pertanahan, Badan Pembinaan BUMD, serta Dinas Sumber Daya Air) menyumbang akurasi yang lebih rendah pada \textit{confusion matrix}. Teknik seperti \textit{focal loss} \citep{zahroFairerPublicComplaint2025}, \textit{class-balanced sampling} yang lebih agresif, atau augmentasi data berbasis sintesis dapat diuji untuk meningkatkan kinerja pada kategori tersebut.

    \item \textbf{Replikasi server dan kebijakan ketersediaan.} Karena server FastAPI menjadi titik inferensi tunggal, kebijakan ketersediaan layanan perlu dirancang dengan cermat. Penelitian lanjutan dapat melakukan replikasi server pada beberapa region, menambahkan \textit{circuit breaker} dan \textit{retry policy} pada klien Flutter, serta merancang prosedur fallback ketika server tidak dapat dijangkau dalam waktu yang lama.

    \item \textbf{Migrasi ke infrastruktur produksi dan optimasi biaya \textit{deployment}.} Server inferensi prototipe yang kini berjalan di HuggingFace Spaces perlu dipindahkan ke infrastruktur produksi terdedikasi, misalnya layanan kontainer terkelola seperti Google Cloud Run, AWS ECS/Fargate, atau Azure Container Apps, yang menyediakan jaminan ketersediaan, \textit{autoscaling}, pemantauan, dan jaringan privat sesuai kebutuhan layanan publik. Sejalan dengan itu, karena inferensi memuat dua \textit{encoder} besar, efisiensi biaya dapat ditingkatkan melalui: (a)~kuantisasi INT8 pada kedua \textit{encoder} sehingga \textit{pipeline} dapat dilayani pada simpul CPU yang jauh lebih murah daripada GPU, mengingat kepala CatBoost sudah ramah-CPU; (b)~penggantian tokenizer berbasis \texttt{transformers} dengan implementasi mandiri (misalnya pustaka \texttt{tokenizers} atau kosakata yang di-\textit{vendor}) agar dependensi \textit{runtime} dan ukuran \textit{image} kontainer berkurang serta permukaan rantai pasok menyempit; (c)~pola penyajian \textit{serverless} atau \textit{batch} dengan \textit{autoscaling} agar biaya proporsional terhadap trafik nyata, bukan biaya simpul yang menyala terus-menerus; serta (d)~audit lisensi \textit{encoder} hulu (DINOv3 dan \textit{multilingual-E5}) untuk memastikan kepatuhan sebelum \textit{deployment} produksi. Studi kelayakan tersendiri dapat membandingkan biaya-akurasi antar-opsi, yaitu server GPU, server CPU terkuantisasi, dan varian \textit{on-device}.
\end{enumerate}
